% ============================================================
% Question Bank: ch08-01 Populations and Samples (Modified — OK/TX)
% Topic Title: Populations and Samples (Experimental Design Emphasis)
% CCSS: 7.SP.A.1 (OK/TX: experimental design emphasis)
% Questions: 27 (15 MC + 6 SA + 6 Graphical)
% ============================================================

% --- Multiple Choice Questions (15) ---

\begin{practiceQuestion}{m-8-1-q01}{mc}
\begin{questionText}
A scientist wants to test whether a new plant fertilizer increases growth. She grows $20$ plants with the fertilizer and $20$ plants without it. The group without fertilizer is called:
\end{questionText}
\choiceA{the sample}
\choiceB{the experimental group}
\choiceC{the control group}
\choiceD{the population}
\correctAnswer{C}
\explanation{The control group does not receive the treatment and is used for comparison. The $20$ plants without fertilizer serve as the control group.}
\end{practiceQuestion}

\begin{practiceQuestion}{m-8-1-q02}{mc}
\begin{questionText}
Which sampling method involves dividing the population into groups (such as grade levels) and then randomly selecting members from each group?
\end{questionText}
\choiceA{Simple random sampling}
\choiceB{Systematic sampling}
\choiceC{Convenience sampling}
\choiceD{Stratified random sampling}
\correctAnswer{D}
\explanation{Stratified random sampling divides the population into subgroups (strata), then randomly samples from each stratum to ensure all groups are represented.}
\end{practiceQuestion}

\begin{practiceQuestion}{m-8-1-q03}{mc}
\begin{questionText}
A researcher wants to determine if a new study method improves test scores. In this experiment, the independent variable is:
\end{questionText}
\choiceA{the test score}
\choiceB{the study method}
\choiceC{the number of students}
\choiceD{the time of the test}
\correctAnswer{B}
\explanation{The independent variable is what the researcher changes or controls — the study method. The test score is the dependent variable because it is measured to see the effect.}
\end{practiceQuestion}

\begin{practiceQuestion}{m-8-1-q04}{mc}
\begin{questionText}
A school has $600$ students. A survey is given to the first $30$ students who arrive in the morning. This sample is most likely:
\end{questionText}
\choiceA{random}
\choiceB{stratified}
\choiceC{biased}
\choiceD{systematic}
\correctAnswer{C}
\explanation{Surveying only early arrivals creates a convenience sample. Students who arrive early may share traits (living close by, riding a bus) that don't represent the full population, making the sample biased.}
\end{practiceQuestion}

\begin{practiceQuestion}{m-8-1-q05}{mc}
\begin{questionText}
In an experiment testing whether music helps concentration, a researcher has one group study with music and another study in silence. Having both groups take the same test at the same time helps to:
\end{questionText}
\choiceA{increase the population size}
\choiceB{control confounding variables}
\choiceC{create a convenience sample}
\choiceD{eliminate the need for a hypothesis}
\correctAnswer{B}
\explanation{Keeping conditions identical (same test, same time) except for the variable being tested (music) controls confounding variables that could affect results.}
\end{practiceQuestion}

\begin{practiceQuestion}{m-8-1-q06}{mc}
\begin{questionText}
A town has $4{,}000$ households. A researcher mails a survey to $200$ randomly chosen households and $120$ respond. The population for this study is:
\end{questionText}
\choiceA{the $120$ who responded}
\choiceB{the $200$ who received the survey}
\choiceC{all $4{,}000$ households}
\choiceD{the researcher}
\correctAnswer{C}
\explanation{The population is the entire group the study is about — all $4{,}000$ households. The $200$ chosen form the sample, and the $120$ who responded are a subset of that sample.}
\end{practiceQuestion}

\begin{practiceQuestion}{m-8-1-q07}{mc}
\begin{questionText}
An experiment is designed to test whether cold water freezes faster than warm water. What is the hypothesis?
\end{questionText}
\choiceA{Water can freeze.}
\choiceB{Cold water freezes faster than warm water.}
\choiceC{Water temperature can be measured.}
\choiceD{An experiment was conducted.}
\correctAnswer{B}
\explanation{A hypothesis is a testable prediction. ``Cold water freezes faster than warm water'' is the specific prediction being tested.}
\end{practiceQuestion}

\begin{practiceQuestion}{m-8-1-q08}{mc}
\begin{questionText}
Every $10$th customer entering a store is asked to complete a survey. This is an example of:
\end{questionText}
\choiceA{convenience sampling}
\choiceB{stratified random sampling}
\choiceC{systematic sampling}
\choiceD{voluntary response sampling}
\correctAnswer{C}
\explanation{Selecting every $k$th individual from a list or sequence is systematic sampling.}
\end{practiceQuestion}

\begin{practiceQuestion}{m-8-1-q09}{mc}
\begin{questionText}
A student wants to find out which lunch option is most popular at school. Which method would produce the most representative sample?
\end{questionText}
\choiceA{Ask all students in the cafeteria line}
\choiceB{Ask $10$ friends}
\choiceC{Randomly select $50$ students from each grade}
\choiceD{Post a survey online and count responses}
\correctAnswer{C}
\explanation{Randomly selecting students from every grade produces a stratified random sample, which best represents the entire school population.}
\end{practiceQuestion}

\begin{practiceQuestion}{m-8-1-q10}{mc}
\begin{questionText}
In a controlled experiment, why is it important to test only one variable at a time?
\end{questionText}
\choiceA{It makes the experiment shorter.}
\choiceB{It ensures you can determine which variable caused the observed effect.}
\choiceC{It eliminates the need for a control group.}
\choiceD{It increases the population size.}
\correctAnswer{B}
\explanation{Testing one variable at a time allows the researcher to attribute any changes in the outcome directly to that variable, making the results valid and interpretable.}
\end{practiceQuestion}

\begin{practiceQuestion}{m-8-1-q11}{mc}
\begin{questionText}
A farmer plants one field with fertilizer and a neighboring field without. The crops grow better in the fertilized field. A classmate says the result may not be valid because the fields could have different soil quality. The classmate is pointing out a potential:
\end{questionText}
\choiceA{hypothesis}
\choiceB{confounding variable}
\choiceC{sample size}
\choiceD{population}
\correctAnswer{B}
\explanation{Soil quality is a confounding variable — a factor other than the treatment (fertilizer) that could influence the outcome, making the result unreliable.}
\end{practiceQuestion}

\begin{practiceQuestion}{m-8-1-q12}{mc}
\begin{questionText}
A survey question reads: ``Don't you agree that school lunches are terrible?'' This question is an example of:
\end{questionText}
\choiceA{a random question}
\choiceB{a fair question}
\choiceC{a leading or biased question}
\choiceD{a stratified question}
\correctAnswer{C}
\explanation{The phrasing ``Don't you agree'' and ``terrible'' push the respondent toward a negative answer, making it a biased or leading question.}
\end{practiceQuestion}

\begin{practiceQuestion}{m-8-1-q13}{mc}
\begin{questionText}
A researcher surveys $40$ students from a school of $800$. What is the sample size?
\end{questionText}
\choiceA{$40$}
\choiceB{$800$}
\choiceC{$760$}
\choiceD{$840$}
\correctAnswer{A}
\explanation{The sample size is the number of individuals actually included in the study, which is $40$.}
\end{practiceQuestion}

\begin{practiceQuestion}{m-8-1-q14}{mc}
\begin{questionText}
Which of the following is the best way to collect data about whether a new recess schedule increases student happiness?
\end{questionText}
\choiceA{Ask one teacher for their opinion}
\choiceB{Survey students before and after the schedule change}
\choiceC{Look up a study from a different school}
\choiceD{Guess based on the length of recesses}
\correctAnswer{B}
\explanation{A before-and-after survey directly collects data from the affected group, allowing comparison and measurement of the schedule's actual effect.}
\end{practiceQuestion}

\begin{practiceQuestion}{m-8-1-q15}{mc}
\begin{questionText}
A study claims that listening to classical music improves math scores. The study tested $5$ students. Why should this conclusion be questioned?
\end{questionText}
\choiceA{Classical music is not real music.}
\choiceB{The sample size is too small to draw reliable conclusions.}
\choiceC{Math scores cannot be measured.}
\choiceD{All studies with music are invalid.}
\correctAnswer{B}
\explanation{A sample of $5$ is too small to produce reliable, generalizable results. Small samples increase the chance that results are due to random variation.}
\end{practiceQuestion}

% --- Short Answer Questions (6) ---

\begin{practiceQuestion}{m-8-1-q16}{sa}
\begin{questionText}
Describe how you would design an experiment to test whether drinking water before a test improves scores. Identify the independent variable, the dependent variable, and the control group.
\end{questionText}
\correctAnswer{Independent variable: whether students drink water before the test. Dependent variable: test scores. Control group: students who do not drink water before the test.}
\explanation{The independent variable is what the experimenter changes (water/no water). The dependent variable is what is measured (test scores). The control group receives no treatment for comparison.}
\end{practiceQuestion}

\begin{practiceQuestion}{m-8-1-q17}{sa}
\begin{questionText}
A school of $500$ students has $200$ in Grade 6, $180$ in Grade 7, and $120$ in Grade 8. A researcher wants a stratified sample of $50$ students. How many should be selected from each grade?
\end{questionText}
\correctAnswer{Grade 6: $20$, Grade 7: $18$, Grade 8: $12$}
\explanation{Proportion from each grade: Grade 6: $\frac{200}{500} \times 50 = 20$. Grade 7: $\frac{180}{500} \times 50 = 18$. Grade 8: $\frac{120}{500} \times 50 = 12$. Total: $20 + 18 + 12 = 50$.}
\end{practiceQuestion}

\begin{practiceQuestion}{m-8-1-q18}{sa}
\begin{questionText}
A company wants to know if customers prefer their new package design. They survey people leaving one store location on a Saturday morning. Give two reasons this sample may be biased.
\end{questionText}
\correctAnswer{(1) Only customers at one location are surveyed, not the full customer base. (2) Only Saturday morning shoppers are included, who may not represent all shoppers.}
\explanation{Both the location and the time of day limit the variety of customers surveyed, leading to a sample that may not represent the entire population.}
\end{practiceQuestion}

\begin{practiceQuestion}{m-8-1-q19}{sa}
\begin{questionText}
Rewrite the following biased survey question as a fair question: ``Why is soccer the best sport to play?''
\end{questionText}
\correctAnswer{A fair version: ``Which sport do you most enjoy playing?''}
\explanation{The original assumes soccer is the best, pushing the respondent. A fair question does not assume a particular answer and allows all options.}
\end{practiceQuestion}

\begin{practiceQuestion}{m-8-1-q20}{sa}
\begin{questionText}
A student claims that since $3$ out of $5$ of her friends like pepperoni pizza, about $60\%$ of all students at school like pepperoni. Explain why this inference may not be valid.
\end{questionText}
\correctAnswer{Her friends are not a random sample — they are a convenience sample. A small, non-random group may not represent the whole school.}
\explanation{Friends likely share similar tastes, so the sample is biased. Also, $5$ people is too small to make reliable inferences about the entire school population.}
\end{practiceQuestion}

\begin{practiceQuestion}{m-8-1-q21}{sa}
\begin{questionText}
Explain the difference between an observational study and an experiment.
\end{questionText}
\correctAnswer{In an observational study, the researcher observes and records data without changing anything. In an experiment, the researcher actively applies a treatment or changes a variable to study its effect.}
\explanation{The key difference is intervention: experiments manipulate a variable while observational studies only observe what naturally occurs.}
\end{practiceQuestion}

% --- Graphical Questions (3 GMC + 3 GSA) ---

\begin{practiceQuestion}{m-8-1-q22}{gmc}
\begin{questionText}
A student surveyed $60$ classmates about their favorite fruit. The bar graph below shows the results.

\begin{center}
\begin{tikzpicture}
\begin{axis}[
    ybar,
    bar width=18pt,
    width=9cm, height=5.5cm,
    ymin=0, ymax=25,
    ylabel={Number of Students},
    symbolic x coords={Apple,Banana,Orange,Grape,Strawberry},
    xtick=data,
    x tick label style={rotate=30, anchor=east},
    nodes near coords,
    every node near coord/.style={font=\scriptsize},
]
\addplot[fill=funBlue!70] coordinates {(Apple,18) (Banana,12) (Orange,15) (Grape,6) (Strawberry,9)};
\end{axis}
\end{tikzpicture}
\end{center}

If there are $300$ students in the school, how many would you predict prefer oranges based on this sample?
\end{questionText}
\choiceA{$15$}
\choiceB{$25$}
\choiceC{$50$}
\choiceD{$75$}
\correctAnswer{D}
\explanation{Orange proportion: $\frac{15}{60} = \frac{1}{4}$. Prediction: $\frac{1}{4} \times 300 = 75$ students.}
\end{practiceQuestion}

\begin{practiceQuestion}{m-8-1-q23}{gmc}
\begin{questionText}
A researcher tested $3$ fertilizers on plant growth. Each fertilizer was used on $10$ plants. The results are shown below.

\begin{center}
\begin{tabular}{|l|c|c|}
\hline
\textbf{Group} & \textbf{Mean Height (cm)} & \textbf{Number of Plants} \\
\hline
Control (no fertilizer) & $12.4$ & $10$ \\
Fertilizer A & $15.8$ & $10$ \\
Fertilizer B & $14.1$ & $10$ \\
Fertilizer C & $13.0$ & $10$ \\
\hline
\end{tabular}
\end{center}

Which fertilizer appears to be the most effective at increasing growth?
\end{questionText}
\choiceA{Control}
\choiceB{Fertilizer A}
\choiceC{Fertilizer B}
\choiceD{Fertilizer C}
\correctAnswer{B}
\explanation{Fertilizer A produced the highest mean height ($15.8$ cm), which is $15.8 - 12.4 = 3.4$ cm taller than the control group. This suggests Fertilizer A is the most effective.}
\end{practiceQuestion}

\begin{practiceQuestion}{m-8-1-q24}{gmc}
\begin{questionText}
The table below shows how students were sampled for a survey.

\begin{center}
\begin{tabular}{|l|c|c|}
\hline
\textbf{Grade} & \textbf{Students in Grade} & \textbf{Sampled} \\
\hline
6 & $180$ & $36$ \\
7 & $150$ & $30$ \\
8 & $120$ & $24$ \\
\hline
Total & $450$ & $90$ \\
\hline
\end{tabular}
\end{center}

What type of sampling method was used?
\end{questionText}
\choiceA{Convenience sampling}
\choiceB{Simple random sampling}
\choiceC{Systematic sampling}
\choiceD{Stratified random sampling}
\correctAnswer{D}
\explanation{Each grade was sampled proportionally ($\frac{36}{180} = \frac{30}{150} = \frac{24}{120} = \frac{1}{5}$). This is stratified random sampling — dividing into strata and sampling each proportionally.}
\end{practiceQuestion}

\begin{practiceQuestion}{m-8-1-q25}{gsa}
\begin{questionText}
A researcher wants to test whether background noise affects reading speed. She has $30$ volunteers.

Describe a controlled experiment she could design. Include:
\begin{itemize}
\item How to divide participants into groups
\item What the independent and dependent variables are
\item What should be kept constant
\end{itemize}
\end{questionText}
\correctAnswer{Randomly divide the $30$ volunteers into two groups of $15$. One group reads in a quiet room (control group) and the other reads with background noise (experimental group). Independent variable: background noise. Dependent variable: reading speed. Keep constant: same passage, same time limit, same room temperature, same lighting.}
\explanation{Random assignment, a control group, and keeping all other conditions the same are essential elements of a well-designed controlled experiment.}
\end{practiceQuestion}

\begin{practiceQuestion}{m-8-1-q26}{gsa}
\begin{questionText}
A school is deciding between three new mascot options. The table shows survey responses from a random sample.

\begin{center}
\begin{tabular}{|l|c|}
\hline
\textbf{Mascot} & \textbf{Votes} \\
\hline
Eagles & $32$ \\
Wolves & $28$ \\
Panthers & $20$ \\
\hline
Total & $80$ \\
\hline
\end{tabular}
\end{center}

The school has $640$ students total. Based on this sample, how many students would you predict would choose each mascot?
\end{questionText}
\correctAnswer{Eagles: $256$, Wolves: $224$, Panthers: $160$}
\explanation{Scale factor: $\frac{640}{80} = 8$. Eagles: $32 \times 8 = 256$. Wolves: $28 \times 8 = 224$. Panthers: $20 \times 8 = 160$. Check: $256 + 224 + 160 = 640$.}
\end{practiceQuestion}

\begin{practiceQuestion}{m-8-1-q27}{gsa}
\begin{questionText}
Two students each conducted a survey to find the favorite subject of $7$th graders at their school.

\textbf{Student A:} Surveyed $10$ students sitting at one lunch table.

\textbf{Student B:} Used a computer program to randomly select $50$ students from the school roster.

Compare the two methods. Which student's results are more likely to represent the whole school? Justify your answer.
\end{questionText}
\correctAnswer{Student B's results are more representative. Student B used a larger, random sample from the whole school. Student A used a small convenience sample that is likely biased because friends at one table may share similar opinions.}
\explanation{Larger random samples better represent the population. Convenience samples (one lunch table) are often biased because they consist of similar individuals. Student B's method of random selection from the full roster is the proper approach.}
\end{practiceQuestion}
